\section{Heroes and Vibhutis}

\begin{quotex}
A happy life is impossible; the highest to which man can attain is an heroic course of life.
\flright{\textsc{Arthur Schopenhauer}}

Then said his Lordship, ``Well. God mend all!''--``Nay, by God, Donald, we must help him to mend it!'' said the other.
\flright{\textsc{Thomas Carlyle}, \textit{Epigraph to Latter-Day Pamphlets}}

The Great Man was always as lightning out of Heaven; the rest of men waited for him like fuel, and then they too would flame.
\flright{\textsc{Thomas Carlyle}}

Suckers try to win arguments, nonsuckers try to win.
\flright{\textsc{Fat Tony}}
\end{quotex}


\paragraph{Avatars}


In \emph{Letters on Yoga}, \textbf{Sri Aurobindo} describes the Avatar as the ``one who comes to open the Way for humanity to a higher consciousness.'' The primary exemplars he mentions are Krishna, Buddha, and Christ, a theme that Valentin Tomberg developed in his quest for the Bodhisattva. Although Avatars have ``the knowledge necessary for their work, they need not have more.'' Specifically, they need not know facts or theories that have no bearing on their mission, and should not be judged by those extraneous standards.

Furthermore, it is not the outer events of their lives that matter most, but rather ``the spiritual Reality, the Power, the Influence that came with him or that he brought down by his action and his existence.'' Hence, the outward events of Christ's life are not as important as the spiritual power that has manifested long after his time. The hylic man is only concerned with the historical and material.

\paragraph{Vibhutis}


The Vibhuti, on the other hand, needs not to be aware of being a power of the Divinity. Nevertheless, he does embody some aspect of the divine, and ``is enabled by it to act with great force in the world.'' When judged by human standards, the Vibhuti may not be recognized as such. Aurobindo explains that

\begin{quotex}
they are not concerned with morality or immorality, perfection or imperfection according to small human standards or setting an example to men or showing new moral attitudes or giving new spiritual teachings. These things may or may not be done, but they are not at all the essence of the matter.
\end{quotex}


Vibhutis may be saints and sage, but they may also be men of action like a Julius Caesar or a Napoleon. There should be no surprises here since Caesar is considered one of the Nine Worthies\footnote{\url{http://www.traditioninaction.org/religious/h097_Worthies.html}}. Aurobindo explains that their faults don't matter:

\begin{quotex}
But do you really believe that men like Napoleon, Caesar ... were not great men and did nothing for the world or for the cosmic purpose? that God was deterred from using them for His purpose because they had defects of character and vices? What a singular idea!
\end{quotex}


\paragraph{The Idea of the Hero}


\textbf{Thomas Carlyle} has written about Vibhutis, at least of a certain type, in \emph{On Heroes, Hero-worship, and the Heroic in History}. This study is of value, Carlyle says, because

\begin{quotex}
Universal History, the history of what man has accomplished in this world, is at bottom the History of the Great Men who have worked here. They were the leaders of men, these great ones; the modelers, patterns, and in a wide sense creators, of whatsoever the general mass of men contrived to do or to attain; all things that we see standing accomplished in the world are properly the outer material result, the practical realization and embodiment, of Thoughts that dwelt in the Great Men sent into the world: the soul of the whole world's history, it may justly be considered, were the history of these.
\end{quotex}


Now some of Carlyle's heroes are not at all appealing to me, and perhaps not to others. For example, I would oppose Oliver Cromwell. Moreover, Cromwell excited a derangement syndrome in his era: his dead body was exhumed, ``executed'', and decapitated, the head being publically displayed for years. But that is not the point. Rather, we should ask, ``did he create history?'' or ``what do we learn about leadership from him?''

Carlyle chooses his heroes from their roles in the world. Oddly, the sequence of roles represents a sort of ``degeneration'' of castes. For example, the first is the Hero as Divinity, exemplified by
\textbf{Odin}, a man-god who articulated a spiritual vision of the world. Next is the Hero as Prophet, no longer a man-god, but rather a messenger of God. This is \textbf{Mohammed}, whom Carlyle regards on a par with the Hebrew prophets of the Old Testament. \textbf{Dante} and \textbf{Shakespeare} are heroes as the inspired Poet. \textbf{Luther} and \textbf{Knox} represent the Hero as Priest; \textbf{Johnson}, \textbf{Rousseau}, and \textbf{Burns}, the Hero as the Man of Letters; and finally, \textbf{Cromwell} and \textbf{Napoleon} as the Hero as King. The word King means the ``one who can''. The Hero begins with a god who reveals the inner nature of the Universe to the Scandinavians, continues with those who are the Prophet, the Poet, the Priest of God, and concludes with those who are active in the world.

\paragraph{Religion}


\begin{quotex}
A man's religion is the chief fact with regard to him. A man's, or a nation of men's.
\end{quotex}


Carlyle does not mean by religion simply a man's nominal affiliation. Incessant debates over creeds do not penetrate deeper than a man's ``argumentative region''. The religious impulse goes much deeper:

\begin{quotex}
The thing a man does practically believe (and this is often enough without asserting it even to himself, much less to others); the thing a man does practically lay to heart, and know for certain, concerning his vital relations to this mysterious Universe, and his duty and destiny there, that is in all cases the primary thing for him, and creatively determines all the rest.
\end{quotex}


That is because:

\begin{quotex}
The thoughts they had were the parents of the actions they did; their feelings were parents of their thoughts: it was the unseen and spiritual in them that determined the outward and actual.
\end{quotex}


So Carlyle tries to get to the soul feeling of a religion, not its professed creed. As examples, he points out the characteristic feelings of various religious impulses:

\begin{itemize}


\item \textbf{Heathenism}: plurality of gods, sensuous representation of the Mystery of Life, physical force as the chief element. A recognition of the forces of Nature as godlike, stupendous, personal Agencies.

\item \textbf{Christianism}: faith in an Invisible, not as real only, but as the only reality. Time, through every meanest moment of it, resting on Eternity; Pagan empire of Force displaced by a nobler supremacy, that of Holiness.

\item \textbf{Scepticism}: uncertainty and inquiry whether there was an Unseen World, any Mystery of Life except a mad one.
\end{itemize}


Such core concepts determine the course of a man's---or a nation's---life. Carlyle dismisses two common theories about the origins of Religion, that are still common today.

\begin{itemize}


\item \textbf{Quackery}. The first is that it is quackery, a contrivance deliberately fabricated to get other men to follow. Perhaps quackery is a factor in the decline of religion, but at its origin, religion had to reach something deep in man in order to take hold.

\item \textbf{Allegory}. The second is that religion is merely allegorical. That would make of religion a game, a sport. But a man's life is not a game, nor will a sane man give up his life for an allegory. A man does not live by a clever allegory, but rather by his deeply held view of the Universe. First there is faith, and only afterward can there be an allegory.
\end{itemize}


Carlyle then turns to Plato's cave. A man escapes from the cave and in astonishment suddenly sees the world anew, the same world that everyone else just takes for granted. Such a man would experience the world as beautiful, awful, and unspeakable. He would experience it directly, and not through the verbal superstructure of words and theories that promise to explain the world to us, but in actuality merely obscure its deep mystery.

Carlyle claims that Odin was such a man who stepped out of the cave. What others may have been feeling, yet unable to articulate, was brought to light by Odin. Carlyle explains:

\begin{quotex}
the great Thinker came, the original man, the Seer; whose shaped spoken Thought awakes the slumbering capability of all into Thought. It is ever the way with the Thinker, the spiritual Hero. What he says, all men were not far from saying, were longing to say.
\end{quotex}


A thought, once awakened, does not slumber again. Odin gave his people the assurance of their destiny, he solved for them the riddle of life. There was also a recognition of the tragic sense of life:

\begin{quotex}
They seem to have seen, these brave old Northmen, what Meditation has taught all men in all ages, That this world is after all but a show---a phenomenon or appearance, no real thing. All deep souls see into that---the Hindoo Mythologist, the German Philosopher, the Shakespeare, the earnest Thinker, wherever he may be: ``We are such stuff as Dreams are made of!''
\end{quotex}


\paragraph{The Character of the Hero}


There is no point in reviewing each of the representative men selected by Carlyle as a Hero. The better course instead is to specify precisely what they represent, i.e., their characters. Carlyle reveals their essences which are easy to miss, given the harsh criticism that is the other side of being a Hero. Sceptics may call such a man false. Yet, how can a falsehood motivate millions of men for decades, centuries, even millennia? If the false is unnatural, then the Hero is \textbf{natural}. Carlyle explains:

\begin{quotex}
A man must conform himself to Nature's laws, be verily in communion with Nature and the truth of things, or Nature will answer him, No, not at all!
\end{quotex}


Moreover, he is \textbf{sincere}, his acts are not calculated:

\begin{quotex}
the Great Man does not boast himself sincere, far from that; perhaps does not ask himself if he is so: I would say rather, his sincerity does not depend on himself; he cannot help being sincere!
\end{quotex}


That leads to the primary definition of the Great Man:

\begin{quotex}
a deep, great, genuine sincerity, is the first characteristic of all men in any way heroic. Not the sincerity that calls itself sincere; ah no, that is a very poor matter indeed;---a shallow braggart conscious sincerity; oftenest self-conceit mainly. The Great Man's sincerity is of the kind he cannot speak of, is not conscious of: nay, I suppose, he is conscious rather of insincerity; for what man can walk accurately by the law of truth for one day? No, the Great Man does not boast himself sincere, far from that; perhaps does not ask himself if he is so: I would say rather, his sincerity does not depend on himself; he cannot help being sincere! The great Fact of Existence is great to him. Fly as he will, he cannot get out of the awful presence of this Reality. His mind is so made; he is great by that, first of all. Fearful and wonderful, real as Life, real as Death, is this Universe to him. Though all men should forget its truth, and walk in a vain show, he cannot. At all moments the Flame image glares in upon him; undeniable, there, there!
\end{quotex}


Thus, the Great Man is \textbf{original}, a messenger from the Infinite Unknown, whether as a god, a prophet, or a poet:

\begin{quotex}
the words he utters are as no other man's words. Direct from the Inner Fact of things---he lives, and has to live, in daily communion with that. ... \emph{It is from the heart of the world that he comes; he is portion of the primal reality of things.} God has made many revelations: but this man too, has not God made him, the latest and newest of all? The ``inspiration of the Almighty giveth him understanding:'' we must listen before all to him.
\end{quotex}


In many cases, his voice cannot be heard, as it is drowned out by the Idols of the Marketplace and the Idols of the Theatre\footnote{\url{http://www.sirbacon.org/links/4idols.htm}}. Lesser men cannot think beyond the known, or better said, beyond what they believe they know.

\paragraph{Faults}


Critics of the Hero typically focus on the man's \textbf{faults}, real or alleged. Lesser men believe that to be the way to truth while in fact it is a barrier to a higher truth. A Great Woman, \textbf{St. Teresa d'Avila}, in \emph{The Way of Perfection} admonishes us:

\begin{quotex}
Look not on our faults.
\end{quotex}


The way to perfection is the overcoming of faults. Perfection is the goal, not the means. The very fear of faults is itself paralyzing. \textbf{Pope Francis} recently said:


\begin{quotex}
Not taking risks, please, no... prudence... Obeying all the commandments, all of them... Yes, it's true, but this paralyzes you too, it makes you forget so many graces received, it takes away memory, it takes away hope, because it doesn't allow you to go forward.
\end{quotex}


Carlyle, too, understands this when he writes:

\begin{quotex}
On the whole, we make too much of faults; the details of the business hide the real centre of it. Faults? The greatest of faults, I should say, is to be conscious of none.
\end{quotex}


Carlyle gives the example of King David as a Hero with faults, yet is the epitome of holiness. Faults are but the outward sign of an inner struggle:

\begin{quotex}
The inner secret of it, the remorse, temptations, true, often-baffled, never-ended struggle of it ... In this wild element of a Life, he has to struggle onwards; now fallen, deep-abased; and ever, with tears, repentance, with bleeding heart, he has to rise again, struggle again still onwards. That his struggle be a faithful unconquerable one: that is the question of questions. We will put up with many sad details, if the soul of it were true.
\end{quotex}


Aurobindo is emphatic on this point. In our fallen world, there will be privations. Since privations, defects, and vices, are non-being, they should not be the focus. He writes:

\begin{quotex}
Why should the Divine care about the vices of great men? Is he a policeman? So long as one is in the ordinary nature, one has qualities and defects, virtues and vices. When one goes beyond, there are no virtues and vices,---for these things do not belong to the Divine Nature.
\end{quotex}


Rather than a defect, the vices of the Vibhutis may be a necessary accompaniment, the shadow side perhaps, of their mission in the world. The reason that the vices of ordinary men are not as noticeable is because their entire lives are unremarkable. Aurobindo expresses it this way:

\begin{quotex}
Men with great capacities or a powerful mind or a powerful vital have very often more glaring defects of character than ordinary men---or at least the defects of the latter do not show so much, being like themselves, smaller in scale.
\end{quotex}


Of course, a low energy person will never become great. Aurobindo concludes:

\begin{quotex}
They great men have more energy and the energy comes out in what men call vices as well as in what men call virtues.
\end{quotex}


\paragraph{The Valet}


\begin{quotex}
No man is a Hero to his valet.
\flright{\textit{French saying}}
\end{quotex}


If that is so, Carlyle asserts that it is not the fault of the Hero, but rather a limitation of the Valet. The latter has a ``valet-soul''.

\begin{quotex}
The Valet does not know a Hero when he sees him!
\end{quotex}


That applies to the Sceptic, also. Carlyle makes this comparison:

\begin{quotex}
The Valet expected purple mantles, gilt sceptres, bodyguards and flourishes of trumpets: the Sceptic of the Eighteenth century looks for regulated respectable Formulas, ``Principles,'' or what else he may call them; a style of speech and conduct which has got to seem ``respectable,'' which can plead for itself in a handsome articulate manner, and gain the suffrages of an enlightened sceptical Eighteenth century!
\end{quotex}


The passing centuries have changed nothing in this regard. The Valet soul of the 21\textsuperscript{st} century is like that of the 18\textsuperscript{th}. They can only be impressed by outward trappings and appearances. If anything, Sceptics today abound more than ever. As a matter of fact, it is a badge of honour to be a Sceptic. The Sceptic rejects anyone who doesn't think, speak, or act in a way the Sceptic considers respectable. Neither the Valet nor the Sceptic can see the obvious:

\begin{quotex}
The sincere alone can recognize sincerity.
\end{quotex}


Unfortunately, a world of Valets and Sceptics will hinder the Hero:

\begin{quotex}
Not a Hero only is needed, but a world fit for him; a world not of Valets;---the Hero comes almost in vain to it otherwise! ... We shall either learn to know a Hero, a true Governor and Captain, somewhat better, when we see him; or else go on to be forever governed by the Unheroic.
\end{quotex}


Aurobindo gives us a similar warning:

\begin{quotex}
People have begun to try to prove that great men were not great, which is a very big mistake. If greatness is not appreciated by men, the world will become mean, small, dull, narrow and tamasic.
\end{quotex}


\paragraph{Summary}


To understand more deeply than the Valet and the Sceptic, it is necessary to be attuned to God's activity in the world, which is the third dimension of History. Thus, human, all-too-human, standards are insufficient. Aurobindo describes it this way:

\begin{quotex}
The Divinity acts according to another consciousness---the consciousness of the Truth above and the Lila below and it acts according to the need of the Lila, not according to men's ideas of what it should or should not do. This is the first thing one must grasp, otherwise one can understand nothing about the manifestation of the Divine.
\end{quotex}


Aurobindo agrees with Carlyle that falsehoods, shams, or quackeries will lack any spiritual impact on the world:

\begin{quotex}
The Divine does not need to suffer or struggle for himself; if he takes on these things it is in order to bear the world-burden and help the world and men; and if the sufferings and struggles are to be of any help, they must be real. A sham or falsehood cannot help. They must be as real as the struggles and sufferings of men themselves
\end{quotex}


Hence, God works through the laws and conditions that he created. Aurobindo puts it this way:

\begin{quotex}
The Divine Consciousness is omnipotent but it has put forth the instrumental personality in Nature, under the conditions of Nature, and it uses it according to the rules of the game ---though also sometimes to change the rules of the game. If Avatarhood is only a flashing miracle, then I have no use for it. If it is a coherent part of the arrangement of the omnipresent Divine in Nature, then I can understand and accept it.
\end{quotex}


Jesus fed the multitudes with the loaves and fishes, but he did not eliminate Poverty. He cured the sick, but did not banish Disease. He raised Lazarus from the dead, yet did not end physical death. If you always need miracles or unusual events to prove the action of God in the world, you will be sure to miss it. Ultimately, we should not look for a material transformation of the world, but rather a spiritual transformation. The subtle rules the dense, and the outer is a manifestation of the inner. Aurobindo concludes:

\begin{quotex}
Whatever is significant in the outward life is so because it is a symbol of what has been realised within himself and one may go on and say that the inner life also is only significant as an expression, a living representation of the movement of the Divinity behind it. ... The teachings of Christ and Buddha are spiritually true not as mere mental teachings but as the expression of spiritual states or happenings in them which by their life on earth they made possible.
\end{quotex}



\flrightit{Posted on 2017-08-08 by Cologero}

\begin{center}* * *\end{center}

\begin{footnotesize}
\begin{sffamily}
\texttt{James on 2017-08-09 at 05:06 said:}

Thanks for another great post.

\hfill

\texttt{Fred on 2017-08-09 at 15:38 said:}

``his dead body was exhumed, ``executed'', and decapitated, the head being publically displayed for years. But that is not the point.''

I wish they did that in France with Voltaire and Rousseau, but instead they let them stay in pantheon with the punishment of listening to mass every day, until a new generation of revolutionaries freed them and honoured them again.

\hfill

\texttt{Matt on 2020-08-08 at 19:21 said:}

I recently came across the Methodist theologian and educator Joseph Wesley Mathews. The following words could very well describe the Vibhuti:

From the book Bending History, he says that history is ``... that movement from the No Longer to the Not Yet. History is not some philosophy a la Marx and Hegel, something you learned in school years ago, which you either flow with or get ground beneath. No, history is that goingonness where human decisions are made. ... wherever you have this goingonness between the No Longer and the Not Yet, there you have the elite cadre in civilization.''

And from a talk he gave:

``We are going to visit the arena of Profound Humanness called ``Integrity''. Sometimes ``integrity'' is reduced to mean a kind of moral uprightness and steadfastness, in the sense of saying, ``He has too much integrity to ever take a bribe.''

But profound integrity goes far beyond this. Sometimes, in order to distinguish it from more limited popular usage, it is called ``secondary integrity''. This is the integrity which is not constrained by limited moralities, however well-intentioned. The integrity that is profound living is the singularity of thrust of a life committed and ordering every dimension of the self towards that commitment. Thus the self is in fact shaped by the self, and focused towards that commitment. You can say that an audacious creation of the self takes place in integrity, without which you are simply the creation of the various forces impacting you in your society.

Thus the basis of integrity is a destinal resolve -- a resolve that chooses and sets your destiny and out of which your whole life is ordered. The object of that resolve is the ultimate decision of each person, and each person makes that choice, consciously or unconsciously. To do so with awareness is the height of man's responsibility. It is incarnate freedom. It is what real freedom looks like.

When man has thus exercised his freedom he realizes that to be true to himself ever thereafter he has a unique position to look at the values of his society. He is no longer bound by the opinions and codes of his fellow-man, but reevaluates then on the basis of their impact on his destinal resolve.

Thus the man of integrity is continuously engaged in a societal transvaluation, a moving across the values of society and reinterpreting them in line with his life's thrust. It does not give him the liberty of ignoring his society, but his obligation transcends the conformity of living within the codes and mores of his society. Thus the man of profound integrity always seems to not quite fit with his fellow-men, but his actions always are appropriate for him, even to those who oppose him.

No matter how odd the man of profound integrity appears to his neighbors, he experiences himself as securely anchored. While he is very clear that this world is not his home, nevertheless he experiences himself as having found his native vale. He experiences an eternal at-one-ness, not so much with the currents and waves of activity around him, but with the deeper trends of history itself. Amid the flux of wavering to and fro that is so evident in others, he experiences an inexplicable rootedness, as though he has sunk a taproot deep into the foundations of the earth itself. Though he experiences his life as a long journey, even an endless journey, towards the object of his resolve, yet he never senses himself as a stranger on the journey It's as if he'd been there before. Original integrity is experienced primarily by this sense of at-one-ness.

Kierkegaard once write a book about this kind of integrity that he titled, ``Purity of Heart is to Will One Thing''. An ancient philosopher focused his wisdom around this integrity with the advice, ``Know yourself, and to your own self, be true.''

\end{sffamily}
\end{footnotesize}

